% =====================================================================
% REPORT_TEMPLATE.tex  --  IEEE two-column paper for AI Academy ML Project
% Compile with:  pdflatex report && pdflatex report
% =====================================================================
\documentclass[conference]{IEEEtran}
% ===== PACKAGES ======================================================
\usepackage[utf8]{inputenc}
\usepackage[T1]{fontenc}
\usepackage{amsmath,amssymb}
\usepackage{graphicx}
\usepackage{booktabs}
\usepackage{array}
\usepackage{hyperref}
\usepackage{listings}
\usepackage{xcolor}
% ===== CODE LISTING STYLE ============================================
\definecolor{codegreen}{rgb}{0,0.55,0}
\definecolor{codegray}{rgb}{0.4,0.4,0.4}
\definecolor{codepurple}{rgb}{0.55,0,0.78}
\definecolor{backcolour}{rgb}{0.97,0.97,0.97}
\lstdefinestyle{pythonstyle}{
  language=Python,
  backgroundcolor=\color{backcolour},
  commentstyle=\color{codegreen},
  keywordstyle=\color{blue},
  numberstyle=\tiny\color{codegray},
  stringstyle=\color{codepurple},
  basicstyle=\ttfamily\footnotesize,
  breakatwhitespace=false,
  breaklines=true,
  keepspaces=true,
  numbers=left,
  numbersep=5pt,
  tabsize=2,
  frame=single,
  rulecolor=\color{black!30}
}
% ===== TITLE & AUTHORS ===============================================
\title{Ensemble Methods: Boosting vs.\ Bagging\\ {\large Final Project Report}}

\author{
\IEEEauthorblockN{Author One, Author Two, Author Three}
\IEEEauthorblockA{AI Academy, National AI Center\\
Spring 2026\\
\{email1, email2, email3\}@example.com}
}

% =====================================================================
\begin{document}
\maketitle

% ----- Abstract -------------------------------------------------------
\begin{abstract}
This paper presents an empirical comparison of two ensemble learning
paradigms---boosting and bagging---using decision‑tree base learners.
We implement a Decision Tree classifier, AdaBoost (with decision
stumps), and a Random Forest from scratch.  Controlled experiments on
real‑world datasets investigate scaling behaviour, noise robustness,
and bias‑variance decomposition.  Unsupervised analysis (PCA,
K‑Means, DBSCAN) supplements the visual exploration of the data.
Our findings show that \emph{[TODO: summarise main result, e.g.,
“AdaBoost outperforms Random Forest on clean data but is more
sensitive to label noise”]}.

\textbf{Deliverables:} In addition to this report, we provide the
complete source code (tagged \texttt{v1.0-final}), a professional
slide deck for the oral defense, and a separate contribution report
detailing each team member's specific work.
\end{abstract}

% ----- Introduction ---------------------------------------------------
\section{Introduction}
\label{sec:intro}
\IEEEPARstart{T}{he} central question of this project is: \textit{``Under
what conditions does boosting outperform bagging, and vice versa, and
why?''}  To answer it, we implement the core algorithms from scratch
and evaluate them on datasets of varying characteristics.

The remainder of the paper is structured as follows.
Section~\ref{sec:related} reviews essential literature.
Section~\ref{sec:methods} describes our implementations.
Section~\ref{sec:experiments} details the experimental setup and
Section~\ref{sec:results} presents the results.
Section~\ref{sec:discussion} discusses the findings, and
Section~\ref{sec:conclusion} concludes.
The accompanying slide deck and contribution report are submitted
separately as required by the project brief.

% ----- Related Work ---------------------------------------------------
\section{Related Work}
\label{sec:related}
\subsection{Decision Trees}
Decision trees recursively partition the feature space to minimise
impurity~\cite{hastie2009elements}.  CART~\cite{quinlan1986induction}
uses Gini impurity or entropy.

\subsection{Boosting}
AdaBoost~\cite{freund1997decision} iteratively re‑weights samples,
giving higher importance to misclassified instances.  The weight
update rule is $\alpha_m = \ln\frac{1-\epsilon_m}{\epsilon_m}$ for
binary classification.

\subsection{Bagging and Random Forests}
Bagging~\cite{breiman1996bagging} reduces variance by training
bootstrap replicates.  Random Forests~\cite{breiman2001random} add
feature sub‑sampling to decorrelate trees further.

\subsection{Unsupervised Techniques}
PCA~\cite{pearson1901lines} finds directions of maximum variance.
K‑Means~\cite{lloyd1982least} partitions data into $k$ clusters,
while DBSCAN~\cite{ester1996density} discovers clusters of arbitrary
shape based on density.

% ----- Methods --------------------------------------------------------
\section{Methods}
\label{sec:methods}
We implemented all algorithms in Python, using only NumPy (and
\texttt{sklearn.metrics} for ARI).  The code is available at
\url{https://github.com/...} (tag \texttt{v1.0‑final}).

\subsection{Decision Tree}
We follow the CART methodology with exhaustive split search.
Both Gini impurity and entropy are supported, with \texttt{sample\_weight}
for the boosting module.

\subsection{AdaBoost}
The discrete SAMME variant~\cite{hastie2009elements} uses
depth‑1 stumps.  Sample weights are updated each round and the
stump is trained on weighted data.

\subsection{Random Forest}
Bootstrap samples are drawn with replacement.  Trees are grown
with feature sub‑sampling (\texttt{sqrt} by default).  OOB error
is computed for each sample using trees where it was not in the
bootstrap.

\subsection{Unsupervised Pipeline}
\begin{itemize}
  \item \textbf{PCA:} Eigen‑decomposition of the covariance matrix.
  \item \textbf{K‑Means:} Lloyd’s algorithm with multiple restarts.
  \item \textbf{DBSCAN:} Brute‑force region queries, expanding
        core points.
\end{itemize}

% ----- Experiments ----------------------------------------------------
\section{Experimental Setup}
\label{sec:experiments}
\subsection{Datasets}
We use three datasets (Table~\ref{tab:datasets}).
\begin{table}[h]
\centering
\caption{Datasets used in the study.}
\label{tab:datasets}
\begin{tabular}{lcc}
\toprule
\textbf{Dataset} & \textbf{Samples} & \textbf{Features}\\
\midrule
Breast Cancer Wisconsin & 569 & 30\\
Adult Income & 48\,842 & 14\\
Covertype (subset) & 5\,000 & 54\\
\bottomrule
\end{tabular}
\end{table}
% TODO: replace with actual datasets used.

\subsection{Preprocessing}
Features were standardised (zero mean, unit variance).  Missing values
were dropped/imputed (justify).  For the imbalanced Covertype subset,
we applied random oversampling of the minority class.

\subsection{Metrics}
We report accuracy, macro F1, and AUC‑ROC (micro‑averaged for
multi‑class).  All experiments use 5‑fold cross‑validation with
a fixed random seed (42).

\subsection{Experiments}
\begin{enumerate}
  \item \textbf{Baseline:} Single tree vs.\ sklearn reference.
  \item \textbf{AdaBoost scaling:} 1--200 stumps.
  \item \textbf{Random Forest scaling:} varying trees and depth.
  \item \textbf{Head‑to‑head:} 100 estimators, 5‑fold CV.
  \item \textbf{Noise robustness:} 5\%, 10\%, 20\% label noise.
  \item \textbf{Bias‑variance decomposition:} 100 bootstrap
        replicates.
  \item \textbf{Unsupervised analysis:} PCA scree, elbow,
        k‑distance plots.
\end{enumerate}

% ----- Results --------------------------------------------------------
\section{Results}
\label{sec:results}
% TODO: replace with actual figures and tables.
\begin{figure}[h]
\centering
\includegraphics[width=\linewidth]{figures/example_scaling.pdf}
\caption{Test accuracy vs.\ number of estimators. (Placeholder)}
\label{fig:scaling}
\end{figure}

Table~\ref{tab:headtohead} summarises the head‑to‑head comparison.
\begin{table}[h]
\centering
\caption{5‑fold CV results (mean $\pm$ std). (Placeholder)}
\label{tab:headtohead}
\begin{tabular}{lccc}
\toprule
\textbf{Model} & \textbf{Accuracy} & \textbf{F1} & \textbf{AUC}\\
\midrule
Single Tree & 0.85$\pm$0.02 & 0.84$\pm$0.03 & 0.91$\pm$0.02\\
AdaBoost & 0.92$\pm$0.01 & 0.91$\pm$0.02 & 0.96$\pm$0.01\\
Random Forest & 0.93$\pm$0.01 & 0.92$\pm$0.01 & 0.97$\pm$0.01\\
\bottomrule
\end{tabular}
\end{table}

The noise robustness experiment (Fig.~\ref{fig:noise}) shows
\emph{[TODO: describe observation]}.  The bias‑variance
decomposition (Table~\ref{tab:biasvar}) confirms that boosting
reduces bias while bagging reduces variance.

\begin{figure}[h]
\centering
\includegraphics[width=\linewidth]{figures/example_noise.pdf}
\caption{Accuracy degradation under label noise. (Placeholder)}
\label{fig:noise}
\end{figure}

\begin{table}[h]
\centering
\caption{Bias‑variance decomposition on Breast Cancer dataset.
(Placeholder)}
\label{tab:biasvar}
\begin{tabular}{lcc}
\toprule
\textbf{Model} & \textbf{Bias$^2$} & \textbf{Variance}\\
\midrule
AdaBoost & 0.05 & 0.10\\
Random Forest & 0.08 & 0.06\\
\bottomrule
\end{tabular}
\end{table}

% ----- Discussion -----------------------------------------------------
\section{Discussion}
\label{sec:discussion}
\emph{[TODO: interpret results, answer the central question, discuss
when boosting outperforms bagging and why.  Relate to dataset
characteristics (noise, dimensionality, class imbalance).  Comment
on unsupervised findings: do cluster structures correlate with
classifier performance?]}

% ----- Conclusion -----------------------------------------------------
\section{Conclusion}
\label{sec:conclusion}
\emph{[TODO: summarise key findings, mention limitations, and suggest
future work (e.g., Gradient Boosting, t‑SNE).]}

% ===== BIBLIOGRAPHY ===================================================
\begin{thebibliography}{00}

\bibitem{hastie2009elements}
T.~Hastie, R.~Tibshirani, and J.~Friedman,
\textit{The Elements of Statistical Learning}, 2nd ed.
Springer, 2009.
[Online]. Available: \url{https://hastie.su.domains/ElemStatLearn/}

\bibitem{quinlan1986induction}
J.~R.~Quinlan,
``Induction of decision trees,''
\textit{Machine Learning}, vol.~1, no.~1, pp.~81--106, 1986.

\bibitem{freund1997decision}
Y.~Freund and R.~E.~Schapire,
``A decision-theoretic generalization of on-line learning and an
application to boosting,''
\textit{J. Comput. Syst. Sci.}, vol.~55, no.~1, pp.~119--139, 1997.

\bibitem{breiman1996bagging}
L.~Breiman,
``Bagging predictors,''
\textit{Machine Learning}, vol.~24, no.~2, pp.~123--140, 1996.

\bibitem{breiman2001random}
L.~Breiman,
``Random forests,''
\textit{Machine Learning}, vol.~45, no.~1, pp.~5--32, 2001.

\bibitem{pearson1901lines}
K.~Pearson,
``On lines and planes of closest fit to systems of points in space,''
\textit{Philosophical Magazine}, vol.~2, no.~11, pp.~559--572, 1901.

\bibitem{lloyd1982least}
S.~P.~Lloyd,
``Least squares quantization in PCM,''
\textit{IEEE Trans. Inf. Theory}, vol.~28, no.~2, pp.~129--137, 1982.

\bibitem{ester1996density}
M.~Ester, H.-P.~Kriegel, J.~Sander, and X.~Xu,
``A density-based algorithm for discovering clusters in large
spatial databases with noise,''
in \textit{Proc. 2nd ACM SIGKDD}, 1996, pp.~226--231.

\end{thebibliography}

% ===== APPENDICES (optional) ==========================================
% \appendices
% \section{Code snippets}
% \lstinputlisting[style=pythonstyle]{../src/trees/decision_tree.py}

\end{document}
